% Zhouyi: Chinese setup for a translated arXiv source, compiled with XeLaTeX.
% zhouyi_tex.py inserts this file right after \documentclass, passes no-math to fontspec on the
% \documentclass line, and comments out the encoding and Type 1 font lines that do not work
% with fontspec (see the font comment below).
% scheme=plain adds CJK support without changing the original template's headings.
\usepackage[UTF8,scheme=plain,fontset=none]{ctex}

% A package may have set global font defaults (XCharter sets Extension=.otf), which would
% turn the name lookups below into file lookups; start from clean defaults.
\defaultfontfeatures{}

% Latin text: keep the template's own fonts. Only when zhouyi_tex.py commented out a Type 1
% font line (times, helvet, courier, mathptmx, tgtermes, tgheros, tgcursor, inconsolata, zlmtt,
% or \rmdefault set to ptm and the like, also inside .cls and .sty files) does it define
% \ZhouyiRm, \ZhouyiSf, \ZhouyiTt, \ZhouyiTtinconsolata or \ZhouyiTtlmmono, and the OpenType
% equivalents are loaded by file name. The specific monospace fonts come last and win.
\ifdefined\ZhouyiRm
  \setmainfont{texgyretermes}[Extension=.otf, UprightFont=*-regular, BoldFont=*-bold,
    ItalicFont=*-italic, BoldItalicFont=*-bolditalic]
\fi
\ifdefined\ZhouyiSf
  \setsansfont{texgyreheros}[Extension=.otf, UprightFont=*-regular, BoldFont=*-bold,
    ItalicFont=*-italic, BoldItalicFont=*-bolditalic]
\fi
\ifdefined\ZhouyiTt
  \setmonofont{texgyrecursor}[Extension=.otf, UprightFont=*-regular, BoldFont=*-bold,
    ItalicFont=*-italic, BoldItalicFont=*-bolditalic]
\fi
\ifdefined\ZhouyiTtinconsolata
  \setmonofont{Inconsolatazi4}[Extension=.otf, UprightFont=*-Regular, BoldFont=*-Bold]
\fi
\ifdefined\ZhouyiTtlmmono
  \setmonofont{lmmono10}[Extension=.otf, UprightFont=*-regular, ItalicFont=*-italic,
    BoldFont=lmmonolt10-bold, BoldItalicFont=lmmonolt10-boldoblique]
\fi

% Chinese text: Noto CJK SC (Source Han), looked up through fontconfig.
% CJK has no italic: emphasis in Chinese uses Kaiti when AR PL UKai is installed,
% otherwise a slanted Song face.
\IfFontExistsTF{AR PL UKai CN}
  {\setCJKmainfont{Noto Serif CJK SC}[BoldFont=Noto Serif CJK SC Bold,
     ItalicFont=AR PL UKai CN, BoldItalicFont=AR PL UKai CN, BoldItalicFeatures={FakeBold=2}]}
  {\setCJKmainfont{Noto Serif CJK SC}[BoldFont=Noto Serif CJK SC Bold, AutoFakeSlant]}
\setCJKsansfont{Noto Sans CJK SC}[BoldFont=Noto Sans CJK SC Bold, AutoFakeSlant]
\setCJKmonofont{Noto Sans Mono CJK SC}

% Chinese captions; the abstract title is bilingual like other headings; References stays English like the entries.
\AtBeginDocument{%
  \renewcommand{\figurename}{图}%
  \renewcommand{\tablename}{表}%
  \renewcommand{\abstractname}{摘要（Abstract）}%
}

% Labels that packages print by themselves: algorithm names and keywords, proofs, contents,
% appendix, cleveref names. begindocument/end runs after every package's and the template's own
% \AtBeginDocument code, so these names win over names set there.
\makeatletter
\newcommand\zhouyi@set[2]{\ifdefined#1\renewcommand#1{#2}\fi}
\AddToHook{begindocument/end}{%
  \@ifpackageloaded{algorithm}{\floatname{algorithm}{算法}}{}%
  \zhouyi@set\listalgorithmname{算法目录}%
  \zhouyi@set\algorithmicrequire{\textbf{输入：}}%
  \zhouyi@set\algorithmicensure{\textbf{输出：}}%
  \zhouyi@set\algorithmicreturn{\textbf{返回}}%
  \zhouyi@set\algorithmicif{\textbf{如果}}%
  \zhouyi@set\algorithmicthen{\textbf{则}}%
  \zhouyi@set\algorithmicelse{\textbf{否则}}%
  \zhouyi@set\algorithmicelsif{\textbf{否则如果}}%
  \zhouyi@set\algorithmicfor{\textbf{对于}}%
  \zhouyi@set\algorithmicforall{\textbf{对所有}}%
  \zhouyi@set\algorithmicdo{\textbf{执行}}%
  \zhouyi@set\algorithmicwhile{\textbf{当}}%
  \zhouyi@set\algorithmicrepeat{\textbf{重复}}%
  \zhouyi@set\algorithmicuntil{\textbf{直到}}%
  \zhouyi@set\algorithmicloop{\textbf{循环}}%
  \zhouyi@set\algorithmicfunction{\textbf{函数}}%
  \zhouyi@set\algorithmicprocedure{\textbf{过程}}%
  \zhouyi@set\algorithmicend{\textbf{结束}}%
  \@ifpackageloaded{algpseudocode}{%
    \algrenewtext{EndFor}{\textbf{结束循环}}\algrenewtext{EndWhile}{\textbf{结束循环}}%
    \algrenewtext{EndLoop}{\textbf{结束循环}}\algrenewtext{EndIf}{\textbf{结束判断}}%
    \algrenewtext{EndFunction}{\textbf{结束函数}}\algrenewtext{EndProcedure}{\textbf{结束过程}}}{}%
  \zhouyi@set\proofname{证明}%
  \zhouyi@set\contentsname{目录（Contents）}%
  \zhouyi@set\appendixname{附录}%
  \@ifpackageloaded{cleveref}{%
    \crefname{figure}{图}{图}\Crefname{figure}{图}{图}%
    \crefname{table}{表}{表}\Crefname{table}{表}{表}%
    \crefname{equation}{式}{式}\Crefname{equation}{式}{式}%
    \crefname{algorithm}{算法}{算法}\Crefname{algorithm}{算法}{算法}%
    \crefname{appendix}{附录}{附录}\Crefname{appendix}{附录}{附录}%
    \crefformat{section}{第~#2#1#3~节}\Crefformat{section}{第~#2#1#3~节}%
    \crefformat{subsection}{第~#2#1#3~节}\Crefformat{subsection}{第~#2#1#3~节}%
    \crefname{theorem}{定理}{定理}\crefname{lemma}{引理}{引理}%
    \crefname{corollary}{推论}{推论}\crefname{proposition}{命题}{命题}%
    \crefname{definition}{定义}{定义}\crefname{remark}{注}{注}\crefname{example}{例}{例}%
    \renewcommand\crefpairconjunction{和}\renewcommand\crefmiddleconjunction{、}%
    \renewcommand\creflastconjunction{和}\renewcommand\crefrangeconjunction{至}%
  }{}%
}
\makeatother
